\chapter{Capabilities y delegación}
\label{cap:capabilities}

POSIX clásico tiene un modelo de control de acceso por usuario / grupo /
otros: cada inode lleva tres tripletes \texttt{rwx}. Para escenarios
multi-tenant o sandboxing fino, ese modelo es insuficiente: no soporta
delegación (``dale a $X$ sólo permiso de lectura sobre este inode, sin
derecho a delegar más''), revocación masiva (``invalide todas las caps
derivadas de $C$''), ni atribución (``¿qué cap autorizó esta lectura?'').

sotFS modela las capabilities como \emph{nodos del TypeGraph}. Cada
operación POSIX se ejecuta bajo un \texttt{CapContext} explícito que
identifica la cap y el dominio del operador. Las invariantes de
monotonicidad las preserva el sistema de reescritura por construcción.

\section{El nodo \NCap}

Recordando el cap.~\ref{cap:typegraph}:

\begin{lstlisting}[language=rust]
pub struct Capability {
    pub id: CapId,
    pub rights: Rights,   // bitmask u8
    pub epoch: u64,       // generación, invalida on revoke
}
\end{lstlisting}

El bitmask \texttt{Rights} ocupa 8 bits
(\citaLine{sotfs-graph/src/types.rs}{150}):

\begin{tabular}{ll}
  \texttt{READ}    & $0x01$ \\
  \texttt{WRITE}   & $0x02$ \\
  \texttt{EXECUTE} & $0x04$ \\
  \texttt{GRANT}   & $0x08$ \\
  \texttt{REVOKE}  & $0x10$ \\
\end{tabular}

\paragraph{Aristas asociadas.}

\begin{itemize}
  \item \egrants{}: \NCap{} $\to$ \NInode. Una cap tiene exactamente una
    arista \egrants{} a un inode (la cap \emph{autoriza} ese inode).
  \item \edelegates{}: \NCap{} $\to$ \NCap. Forma el árbol CDT
    (\emph{capability derivation tree}): cada cap (excepto las raíz)
    tiene un padre.
\end{itemize}

\section{El invariante central: monotonicidad}

\begin{invariante}[Cap monotonicity, G4]
\label{inv:cap-mono}
Para todo par $c, c' \in \NCap$ con arista \edelegates{} de $c$ a $c'$,
\[
  c'.\texttt{rights} \;\subseteq\; c.\texttt{rights}.
\]
Una delegación nunca crea derechos; sólo restringe los del padre.
\end{invariante}

\section{Acciones formales en el spec TLA}

\citaTLA{sotfs\_capabilities.tla} (291 LOC) modela cuatro acciones:

\begin{description}
  \item[\textsc{Derive}.] Aloca una cap nueva con $\texttt{rights}
    \subseteq \mathit{parent}.\texttt{rights}$, crea arista
    \edelegates{} desde el padre. Nuevo \texttt{epoch} unique.
  \item[\textsc{Revoke}.] Remueve una cap y \emph{recursivamente todo
    su subárbol} \edelegates{}. Es la acción que justifica el
    \texttt{epoch}: stale handles a la cap revocada se detectan al
    chequear que el \texttt{epoch} corresponde a una cap viva.
  \item[\textsc{Grant}.] Transfiere una cap entre dominios. No afecta
    \texttt{rights} ni el árbol; sólo cambia la atribución.
  \item[\textsc{Access}.] Registra una operación
    \texttt{(cap, inode, right)} en un log de auditoría.
\end{description}

El invariante TLA correspondiente:

\begin{lstlisting}[language=tla]
MonotonicAttenuation ==
    \A c \in caps :
        delegatesParent[c] # NULL =>
            (delegatesParent[c] \in caps =>
                capRights[c] \subseteq capRights[delegatesParent[c]])
\end{lstlisting}

TLC verifica este invariante exhaustivamente bajo configuración
\texttt{small} (3 caps, 3 derives, 1 revoke).

\section{El \texttt{CapContext}: cap-mediated context}

Cada operación POSIX en \citaCrate{sotfs-ops} se ejecuta bajo un
contexto que dice ``qué cap autorizó esto y qué dominio lo emitió''.
La estructura (\citaLine{sotfs-graph/src/types.rs}{189}):

\begin{lstlisting}[language=rust]
pub struct CapContext {
    pub cap_id: Option<CapId>,
    pub domain_id: u64,
}
\end{lstlisting}

\texttt{cap\_id = None} es el path \emph{anonymous} (no se presentó cap;
e.g.\ tarea de admin interna); \texttt{domain\_id = 0} es el dominio
\emph{kernel} privilegiado.

\paragraph{Cómo lo setean los callers.}

\begin{itemize}
  \item \fuse{} callbacks: derivan \texttt{CapContext} del \texttt{Request}
    de FUSE (uid / gid del proceso que invocó la syscall). Ver
    \citaLine{sotfs-fuse/src/fs.rs}{59}.
  \item Tests: setean explícitamente con \texttt{g.set\_cap\_ctx(\dots)}.
  \item \texttt{sotfsctl}: corre con \texttt{CapContext::anonymous()}.
\end{itemize}

\paragraph{El plumbing de v0.2.5} (\citaCommit{61e7faf}). Pre-v0.2.5 el
runtime exponía la API de caps pero \emph{descartaba el contexto en
tiempo de operación}. Cada \texttt{record\_prov} guardaba
\texttt{(cap=None, domain=0)} aunque la regla DPO se hubiera
ejecutado bajo una cap real. El commit cierra ese gap:
\texttt{set\_cap\_ctx} se llama al inicio de cada callback FUSE,
\texttt{record\_prov} lee el contexto vivo y lo serializa al sidecar
JSONL (cap.~\ref{cap:fuse}).

\section{Algoritmo de chequeo de acceso}

\begin{algorithm}[H]
\SetAlgoLined
\KwIn{Grafo $g$, cap $c$, inode objetivo $i$, derecho $r$ requerido}
\KwOut{\textbf{ok} si la cap autoriza la op, \textbf{fail} si no}
\caption{\textsc{CheckAccess}}\label{alg:check-access}

\If{$c \notin g.\mathrm{caps}$}{\Return{\textbf{fail (cap no existe)}}\;}
\If{$r \not\in c.\texttt{rights}$}{\Return{\textbf{fail (rights insuficientes)}}\;}
$\mathit{tgt} \gets g.\mathrm{cap\_grants}[c]$\;
\If{$\mathit{tgt} \neq i$}{\Return{\textbf{fail (cap no autoriza este inode)}}\;}
\Return{\textbf{ok}}\;
\end{algorithm}

\paragraph{Costo.} $O(\log |\NCap|)$ — todo es lookup en BTreeMap. La
complejidad real está en la cadena \edelegates: para verificar que la
cap proviene de un origen confiable, hay que recorrer al padre hasta una
raíz, lo cual es O(profundidad del CDT) en peor caso. En la práctica
las cadenas son cortas (3-4 niveles).

\section{Revocación y el rol del \texttt{epoch}}

Cuando \textsc{Revoke} se ejecuta sobre una cap $c$:

\begin{enumerate}
  \item Calcular el subárbol $\mathrm{Desc}(c)$ siguiendo \edelegates.
  \item Para cada $c' \in \{c\} \cup \mathrm{Desc}(c)$, remover el nodo
    de la arena de caps y todas sus aristas adyacentes.
  \item Las \emph{handles externas} a esas caps (e.g.\ un cliente que
    cacheó un \texttt{CapId}) detectan la revocación al hacer un check
    de \texttt{epoch}: el \texttt{epoch} corriente del cap activo es
    distinto del cacheado, o la cap no existe.
\end{enumerate}

\begin{lema}[Revocación es transitiva]
\label{lem:revoke-transitivo}
Si $c \in \NCap$ es revocada, ninguna cap derivada de $c$ (vía
\edelegates) sobrevive en el grafo. Toda invocación posterior con
cualquier $c' \in \mathrm{Desc}(c)$ falla en \textsc{CheckAccess} en su
primera línea.
\end{lema}

\section*{Resumen del capítulo}

\begin{itemize}
  \item Las capabilities son nodos \NCap{} con \texttt{rights} (bitmask
    u8), \texttt{epoch} y aristas \egrants{} (a \NInode) y \edelegates{}
    (a otra \NCap).
  \item Invariante monotónica: \texttt{rights} de hijos delegados es
    subset del padre. Verificada en \citaLine{sotfs-graph/src/graph.rs}{1056}
    y en \citaTLA{sotfs\_capabilities.tla}.
  \item Cuatro acciones: \textsc{Derive}, \textsc{Revoke},
    \textsc{Grant}, \textsc{Access}.
  \item \texttt{CapContext} se plumbea desde \fuse{} a través de
    \texttt{set\_cap\_ctx} antes de cada op (\citaCommit{61e7faf}); el
    provenance log atribuye cada op a su cap y dominio.
  \item Revocación es transitiva: borrar una cap remueve todo su
    subárbol \edelegates.
\end{itemize}

\section*{Ejercicios}

\begin{ejercicio}[\dificultad{$\star$}]
\label{ej:cap-1}
Liste los cinco bits del \texttt{Rights} bitmask de sotFS.
\end{ejercicio}

\begin{ejercicio}[\dificultad{$\star$}]
\label{ej:cap-2}
¿Cuál es la diferencia entre \texttt{cap\_id = None} y \texttt{cap\_id
= Some(0)}? Pista: revisar \texttt{CapContext::anonymous()}.
\end{ejercicio}

\begin{ejercicio}[\dificultad{$\star\star$}]
\label{ej:cap-3}
Demuestre el \lemref{lem:revoke-transitivo} usando inducción sobre la
profundidad del subárbol \edelegates.
\end{ejercicio}

\begin{ejercicio}[\dificultad{$\star\star$}]
\label{ej:cap-4}
Considere un atacante que obtuvo una \texttt{CapId} cacheada antes de
una revocación. Mostrá cómo el campo \texttt{epoch} hace inútil ese
handle.
\end{ejercicio}

\begin{ejercicio}[\dificultad{$\star\star$}]
\label{ej:cap-5}
Diseñe la prueba TLA correspondiente a la \invref{inv:cap-mono} (sin
mirar \citaTLA{sotfs\_capabilities.tla}). ¿Cuál es la propiedad de
inducción y cuál el estado base?
\end{ejercicio}

\begin{ejercicio}[\dificultad{$\star\star\star$}]
\label{ej:cap-6}
Extienda el modelo a \emph{revocación temporal}: una cap puede
suspenderse por un intervalo $[t_0, t_1]$ y reactivarse después.
Discutí qué nuevos invariantes hay que agregar y qué TLA spec hay que
cambiar.
\end{ejercicio}
